\subsection{Puertas de versión y secuencia de recepción}
La disponibilidad de una versión para migración se distingue de la certificación de la solución completa. Se establecen las puertas VM-E1 y VM-E2 para los ensayos de datos; VS-E1 y VS-E2 para pruebas integrales; y las recepciones H5/H10 y H7/H12 para certificación y producción. Esta separación precisa el alcance de cada evidencia y no reduce los dos ensayos completos exigidos por RT-05.13 ni sustituye la aceptación contractual \parencite[RT-05.13]{bases-transversales}\parencite[art.~18; Formulario~E-25]{bases-admin}.

\begin{longtable}{L{21mm}L{23mm}L{107mm}}
\caption{Puertas de disponibilidad y dependencias}\label{tab:sd7-puertas-version}\\\hline
Puerta & Objetivo relativo & Condición de liberación\\\hline\endfirsthead
\hline Puerta & Objetivo relativo & Condición de liberación\\\hline\endhead
VM-E1 & Antes del primer ensayo en 10 & PREPROD recibido; modelos y contratos de todo el universo a migrar E1; cargadores, archivo, delta y retorno; conciliación reproducible y mecanismos funcionales de verificación de datos. La programación actual ocupa hasta el cierre de 9 y debe anticiparse para absorber incertidumbre.\\\hline
VM-E2 & Antes del primer ensayo en 16 & Modelos y contratos del universo E2, incluidas extensiones de innovaciones que afecten sus datos; transformaciones y conciliación; retorno compatible con E1 en producción. No basta con una entrega parcial de lógica en 15.\\\hline
VS-E1 & Antes de QA y pruebas integrales E1 & Núcleos E1, servicios comunes, canales e interfaces integrados y versionados; protocolo integral, datos y recursos de prueba disponibles. Su liberación se acredita por sus paquetes terminados, sin usar VM-E1 como sustituto.\\\hline
VS-E2 & Antes de QA y pruebas integrales E2 & Núcleos E2, adaptador contable, canales y funcionalidades de innovación aplicables a certificación integrados. Vistas previstas en 16 e interfaces de innovación en 17 impiden presumir esta puerta al cierre de 15.\\\hline
H5/H10 & Meses 12/18 & Pruebas integrales y UAT recibidas sobre versión identificada, ensayos completos válidos, observaciones subsanadas y revisión contractual concluida. Las pruebas propias de innovación de 18 requieren integración y revisión dentro de H10.\\\hline
H7/H12 & Meses 16/21 & Marcha blanca concluida conforme a las bases, capacitación y presencia disponibles, corte autorizado fuera de congelamientos, conciliación y reversión válidas para la versión que se activa.\\\hline
\end{longtable}
\FuenteTabla{Puertas internas de Synaptix; hitos obligatorios de Bases Administrativas E-25 y controles de migración de Bases Técnicas Transversales.}

Las dependencias \texttt{VERSION-E1} y \texttt{VERSION-E2} de los registros de pruebas y capacitación corresponden a VS-E1 y VS-E2, respectivamente. Los nodos E y L de la subred CPM corresponden a VM-E1 y VM-E2. Arquitectura/Integración registrará para cada puerta el manifiesto de componentes, esquemas, contratos, configuración y versiones. Datos firmará la correspondencia entre universo, transformación y conciliación; Calidad registrará qué pruebas cubren ese manifiesto. El responsable de Proyecto comprobará fechas de entrega, revisión y subsanación. Se incorpora al seguimiento de los paquetes existentes; cualquier trabajo adicional se estimará y asignará sin presumir que está cubierto por la reserva.

\paragraph{Tratamiento de cambios posteriores al ensayo.}
Toda modificación de esquema, interpretación de campos, autoridad, interfaz de escritura, conciliación o retorno obliga a analizar su impacto en los ensayos. Si invalida sus resultados, se repetirá el ensayo completo afectado sobre la nueva versión y se conservarán al menos dos ensayos completos válidos antes de la migración definitiva. Un cambio de vista sin impacto demostrado en esos elementos puede conservar el ensayo de datos, pero requiere su prueba funcional y regresión. Calidad documentará la decisión; no se presume compatibilidad por el nombre del componente. No se acortan las revisiones de diez días hábiles ni se desplaza la migración a una ventana protegida para absorber la repetición.

\paragraph{Orden de pruebas sobre ambientes independientes.}
Carga/estrés, continuidad desconectada, resiliencia y conmutación real no se ejecutarán simultáneamente sobre el mismo ambiente si una altera la condición de otra. Arquitectura reservará instancias aisladas o las programará en secuencia, con restauración de línea base entre pruebas. La prueba ofensiva y sus correcciones preceden a la regresión final y a UAT. UAT recibe los resultados de las demás pruebas y ensayos, sin validarse con una versión distinta ni tratar la disponibilidad de refuerzos como disponibilidad automática del ambiente o de la contraparte.

La secuencia corregida conserva 25 enlaces resueltos entre paquetes de una misma quincena y los ordena mediante tiempos de ejecución directa. Las dependencias externas y las puertas simbólicas conservan responsable, fecha requerida y evidencia; no se eliminan de la ruta por carecer todavía de un código de actividad.

El registro \path{dependencias_paquetes.csv} identifica 453 enlaces resueltos entre los 813 paquetes fijos: 400 unen quincenas anteriores, 25 comparten quincena y 28 dependen de meses estacionales pendientes. No hay ciclos en este conjunto resuelto; esto no acredita la red completa mientras falten sus condiciones. \path{condiciones_programacion.csv} conserva 377 referencias a condiciones sin actividad instanciada, con responsable de paquete y fecha requerida; siete corresponden a la familia variable \texttt{IN2.5.4.*}. Esas revisiones previas a botaduras se materializarán por lotes y se enlazarán sólo con sus embarcaciones y períodos aplicables; no se exige esperar a botaduras futuras para cerrar un registro anterior. Su esfuerzo se conserva mediante la fórmula de T-14 y se agregará a T-15 antes de programarlas.

\paragraph{Secuencia corregida y recursos de ejecución directa.}
Synaptix sustituye el reparto que dejaba once paquetes fuera de quincena por una secuencia que traslada esos consumidores a la segunda quincena de su mismo mes. Los servicios comunes y controles se ejecutan primero; la lógica funcional espera su liberación y termina antes de las vistas de 9. Notificaciones espera a consentimiento. Se conservan los códigos, meses, entregables, responsables y HH de cada paquete. La Tabla~\ref{tab:sd7-cambios-quincena} identifica las modificaciones; el diccionario y las actividades correspondientes adoptan esas ventanas.

\begin{longtable}{L{31mm}R{18mm}R{26mm}R{26mm}}
\caption{Corrección de quincenas sin alterar meses ni HH}\label{tab:sd7-cambios-quincena}\\\hline
Paquete & Mes & Quincena anterior & Quincena propuesta\\\hline\endfirsthead
\hline Paquete & Mes & Quincena anterior & Quincena propuesta\\\hline\endhead
1.4.1.1.2 & 8 & 1 & 2\\\hline
1.4.1.3.2 & 8 & 1 & 2\\\hline
1.4.2.1.2 & 8 & 1 & 2\\\hline
1.4.2.3.2 & 8 & 1 & 2\\\hline
1.4.3.2.2 & 8 & 1 & 2\\\hline
1.5.1.1.2 & 8 & 1 & 2\\\hline
1.5.2.2.2 & 8 & 1 & 2\\\hline
1.5.3.2.2 & 8 & 1 & 2\\\hline
1.6.1.2.2 & 8 & 1 & 2\\\hline
1.7.3.1 & 8 & 1 & 2\\\hline
1.7.4.1 & 9 & 1 & 2\\\hline
\end{longtable}\FuenteTabla{Reprogramación de Synaptix por precedencias de servicios comunes, controles y consumidores. Esfuerzo de los paquetes conservado.}


El registro \path{secuencia_fija_recursos.csv} programa los 797 paquetes con mes asignado en una escala de diez días hábiles por quincena. Cada contribución no nula de perfil tiene un ejecutor identificado en \path{asignacion_fija_recursos.csv}; las funciones de una misma persona se suman. Se mantiene un máximo de 50 HH directas por quincena y 100 por mes, con máximo diario de 7,2 HH, equivalente a la referencia de 144 HH productivas en veinte días. Este límite diario permite concentrar esfuerzo en parte de la quincena sin aumentar su capacidad total ni considerar la reserva íntegra disponible en una jornada ocupada.

Para desarrollar esta secuencia, las contribuciones de ejecución se realizan en paralelo conforme a las actividades del paquete, y Calidad revisa después de su finalización. Con $h_{k,p}$ horas por perfil y $v=7{,}2$ HH/día, la duración directa adoptada es
\[
d_k=\max_{p\ne Q}\left(\frac{h_{k,p}}{v}\right)+\frac{h_{k,Q}}{v}.
\]
Se omite el primer término cuando sólo interviene Calidad. Las contribuciones de una misma persona no pueden superar el límite diario conjunto; si coinciden funciones O/L, se asigna refuerzo para evitar esa superposición. La duración expresa ejecución y revisión interna, sin añadir implícitamente recepción de la contraparte ni intervalos completos de observación. No se divide una contribución funcional entre personas para reducir artificialmente su duración.

La asignación busca reutilizar los puestos compatibles con la ventana y su capacidad, y agrega puestos especializados cuando no puede hacerlo. No acredita dotación mínima. Los resultados conservan 20.820 HH en 2.680 contribuciones, comprueban los 425 enlaces fin a inicio entre paquetes con mes fijo y mantienen cada contribución dentro de su quincena, sin sobreasignación diaria en el modelo. La dotación mensual resultante se presenta en T-15 y sustituye la del reparto anterior. Los 28 enlaces estacionales restantes mantienen fecha por definir.

\paragraph{Límites temporales que deben incorporarse a la red.}
La secuencia directa no reduce una observación quincenal a las horas necesarias para redactar su registro. Los cierres de datos de pilotos y los informes que dependen de ellos requieren sus ventanas reales y revisión final posterior al cierre; la red completa incorporará ese evento temporal y recalculará cualquier contribución que deba cambiar de quincena. También se incorporarán recepción y subsanación contractual, ciclos completos de ensayo, capacitación por turnos, reservas, disponibilidad del origen, ambientes y fechas externas. Las 328 HH estacionales y trabajos variables deben recibir sus propios recursos.

Los enlaces a VM/VS y las demás condiciones de liberación conservan la evidencia exigida. La programación directa no se presenta como cumplimiento de esas puertas ni de los hitos hasta integrarlas con el calendario protegido. Proyecto consolidará esa red, comprobará el paralelismo de las actividades y actualizará dotación, equipos y costos antes de aprobar la línea base. La suficiencia de la oferta se evalúa con todos esos elementos, no sólo con el balance de trabajo fijo.


\subsection{Observación, presentación y recepción de pilotos}
La ejecución directa del registro se distingue del período observado. IN1 admite solicitudes ingresadas durante los primeros 21 días naturales del mes y las sigue durante siete días desde su envío; el último seguimiento termina el día 28. Se conservan las solicitudes sin respuesta y las condiciones de comparabilidad. IN4 mide el tiempo activo de las solicitudes de ese mismo universo, con indicadores y exclusiones propios. IN5 conserva 28 días naturales de base y otros 28 de evaluación, con la acción y sus evidencias entre ambos. Los paquetes quincenales realizan captura y control durante la ventana; no acreditan su cierre definitivo al consumir sus HH antes del último dato.

El registro \path{ventanas_observacion.csv} sitúa las bases en 1--28 de junio de 2028 y las evaluaciones en 1--28 de julio de 2028, para el escenario de inicio de diciembre de 2026 y la programación de origen de SD13. Las acciones IN5 se ubican en 29--30 de junio, con acceso permitido y habilitación validada. La formación IN1 depende del cierre de su base y debe terminar antes de activar la cohorte del piloto. Sus sesiones y evaluación requieren una asignación específica en los días restantes; las HH del paquete no prueban por sí solas que esa ventana sea suficiente.

Calidad emitirá los informes definitivos después del cierre de los datos, conservando versiones, participantes, exclusiones, casos pendientes e indicadores. La contraparte dispondrá de diez días hábiles de revisión y, si formula observaciones, Synaptix dispondrá de diez días hábiles para subsanarlas. La subsanación requiere carga y ejecutores identificados; una ventana temporal no aporta HH de corrección implícitas. La repetición de observaciones de la misma naturaleza mantiene el tratamiento de atraso previsto en las bases \parencite[art.~18.2--18.4; p.~13]{bases-admin}.

\begingroup\small\setlength{\tabcolsep}{3pt}
\begin{longtable}{L{41mm}L{24mm}L{24mm}L{24mm}L{24mm}}
\caption{Recepción de pilotos y margen hasta H12 en el escenario de inicio diciembre de 2026}\label{tab:sd7-recepcion-pilotos}\\\hline
Escenario & Informe & Fin revisión & Fin subsanación & Transferencia\\\hline\endfirsthead
\hline Escenario & Informe & Fin revisión & Fin subsanación & Transferencia\\\hline\endhead
Reparto directo secuencial & 03/08/2028 & 18/08/2028 & 01/09/2028 & 07/09/2028\\\hline
Preparación previa y cierre paralelo de informes & 30/07/2028 & 11/08/2028 & 28/08/2028 & 30/08/2028\\\hline
Anticipación de base y piloto a meses 18/19 & 04/07/2028 & 18/07/2028 & 01/08/2028 & 03/08/2028\\\hline
\end{longtable}\FuenteTabla{Escenarios de Synaptix: cierre de observación 28/07/2028 en las dos primeras filas y 28/06/2028 en el anticipo; revisión y subsanación diez días hábiles cada una; lunes a viernes, con 15 de agosto bloqueado por hipótesis conservadora. H12 se contrasta con 31/08/2028; se requiere validar feriados y calendario efectivo.}\endgroup


La subred \path{red_observacion_recepcion.csv} conserva los eventos de cierre, presentación, revisión, subsanación y validación final de transferencia. El primer escenario adopta cuatro días hábiles para terminar un informe de 32 HH conforme a la duración directa de T-15, después del cierre de observación. Con la reserva completa de subsanación y transferencia secuencial, termina el 7 de septiembre, fuera de H12. La programación no puede considerar el informe concluido en una quincena anterior a su último dato.

El segundo escenario estudia preparación previa de insumos y dos frentes de cierre del informe durante el 29--30 de julio. Conserva sus 32 HH y reparte la contribución de Calidad de 16 HH en dos revisiones de 8 HH, sobre secciones identificadas y coordinadas por el responsable de Calidad. El trabajo de fin de semana requiere turnos, descansos y disponibilidad reservados; no se atribuye a una autorización o ejecución realizadas. La documentación de transferencia se prepara durante la revisión y la validación final posterior conserva ocho HH de Calidad en dos días hábiles. Una corrección que invalide esa preparación obliga a recalcular su esfuerzo y duración. Este escenario no sustituye la asignación activa de T-15 hasta incorporar esos ejecutores e intervalos.

Incluso ese cierre paralelo deja sólo el 31 de agosto antes de H12 después de agotar revisión y subsanación. No absorbe los diez días hábiles de reserva marina planteados por campaña. La corrección documental y las cancelaciones por marejada no se imputan a una misma reserva como si nunca pudieran coincidir. El cálculo conserva lunes--viernes y bloquea el 15 de agosto por hipótesis conservadora; se completará con feriados y ventanas efectivas. Ninguno de estos dos escenarios acredita suficiencia integral.

\paragraph{Objetivo de anticipación y control de la línea base.}
El recorrido hacia atrás desde el 31 de agosto de 2028, reservando diez días hábiles de revisión, diez de subsanación, dos de validación final y diez de reserva marina, exige presentar el expediente a más tardar el 17 de julio en este escenario. El cierre de observación del 28 de julio impide cumplir ese objetivo con la programación de origen, aunque se agreguen revisores. Proyecto reprogramará la cadena completa de preparación, construcción, pruebas, base y piloto de las innovaciones afectadas, conservando sus duraciones y condiciones de evaluación, y reconciliará las fechas con SD13, T-14, T-15 y T-18 antes de aprobarla. Anticipar sólo la fecha del informe sin adelantar sus datos no es una solución admisible.

La anticipación se desarrolla trasladando base y evaluación a los meses 18 y 19 y adelantando un mes los paquetes de innovación de 14--20, mientras la materialización contractual conserva el mes 21. El ensayo de programación conserva las 20.820 HH, los 425 enlaces resueltos entre paquetes fijos y el pico de 35 refuerzos en el mes 8. No acredita las puertas simbólicas, las habilitaciones ni el calendario efectivo. Se conserva como alternativa reproducible en \path{ventanas_anticipo_propuesto.csv}, pendiente de incorporar los eventos reales de formación y cierre y reconciliar el calendario de SD13.

En esta alternativa las bases cierran el 28 de mayo, la evaluación el 28 de junio y el informe se presenta el 4 de julio bajo las hipótesis del cálculo. Revisión y subsanación concluyen el 1 de agosto; la validación final termina el 3 de agosto. Los diez días hábiles de reserva marina terminan el 18 de agosto, antes de H12. Este margen temporal permite continuar el diseño sin acortar la observación ni las revisiones. Requiere verificar que la construcción y las pruebas anticipadas estén liberadas y que la capacitación posterior a la base quepa en los turnos previos al 1 de junio; no se atribuye aceptación por consumir ese margen. Hasta completar esa comprobación, la asignación activa conserva sus resultados de esfuerzo y queda condicionada por estos eventos temporales.

IN3 fijará antes de medir la ventana de muestra manual/asistida, su último caso y el cierre de revisión humana; su informe dependerá de esos eventos y de su recepción, sin atribuirle automáticamente las ventanas de IN1 o IN5. IN2 mantiene sus cierres estacionales, revisiones por lotes antes de botaduras y aceptación propia. Los informes de innovación no sustituyen la verificación del volumen real durante las cuatro semanas finales de marcha blanca obligatoria ni sus criterios de aceptación \parencite[art.~17.3; pp.~12--13]{bases-admin}.

